\chapter{Legal, Social, Ethical and Professional Issues}
\label{ch:lsesp}

This chapter is the first place in the dissertation where the project is examined as
a piece of professional practice rather than as a piece of computer science. Three
properties of the work give it a distinctive ethical surface: the corpus consists of
personal-data instruments under EU and UK data-protection law (\S\ref{sec:lsesp-legal}),
the system being studied is general-purpose enough that it could be deployed for
purposes far beyond the academic question (\S\ref{sec:lsesp-social}), and the experiment
itself produced sustainability-relevant measurements that bear on how the technology
should be deployed (\S\ref{sec:lsesp-environment}). The British Computer Society
Code of Conduct~\cite{bcs2022} structures the discussion: each of its four areas is
treated against the specific obligations the project incurs, rather than restated in
the abstract.

\section{Legal context}
\label{sec:lsesp-legal}

\subsection{GDPR and the personal-data character of utility bills}

A utility bill names the customer, the service address, and the account, and ties
those identifiers to a consumption history. It is therefore personal data within
the meaning of GDPR Article 4(1)~\cite{eu_gdpr_2016}, and most of the corpus would
also fall under the UK GDPR retained equivalent. Article 4(1) is not satisfied by
removing the customer's name alone: the address and the account number are themselves
identifiers, and the consumption history is identifiable when joined to a bill. Any
research use of utility bills inherits this status from the document class itself,
not from any individual case.

Three of the principles in Article 5 bear directly on the experimental setup.
Lawfulness and purpose limitation in Article 5(1)(a)--(b) require that the data be
processed for specified, explicit purposes; in this study, the only purpose is the
modality and provider comparison documented in Chapters~\ref{ch:objectives}--\ref{ch:evaluation},
and bills are not retained for any other use. Data minimisation in Article 5(1)(c)
requires that no more data be processed than is necessary; the schema in
\S\ref{sec:schema} restricts extraction to twelve fields, and the per-document
artefacts written in \S\ref{sec:exec-repro} record only the prediction, the
ground truth, the LLM output, and the OCR text, none of which carries
identifiers beyond what is unavoidable for the comparison. Storage limitation in
Article 5(1)(e) is implemented by storing the corpus inside the project repository
under \texttt{data/bills/} as redacted PDFs, with personally identifiable account,
customer and address detail removed before any commit, and by treating the redacted
form as the canonical record.

The lawful basis for processing is Article 6(1)(a) (consent) for the Stream B
documents contributed by personal contacts (\S\ref{sec:dataset-provenance}), and
Article 6(1)(f) (legitimate interests) for the Stream A documents released
publicly by the issuers themselves. The legitimate-interests balancing test for
Stream A is straightforward: the bills are already in the public domain on
issuer help pages and tariff-board portals, redaction further reduces the residual
identifiability, and the comparative-analysis purpose is one a reasonable data
subject would not object to. For Stream B, written consent was obtained before any
contributor's bill entered the corpus, and the consent form covered scope, storage
location, and the right to withdraw.

The single legitimate-interests case worth flagging is the Italian subset, where the
manifest records two bills from \texttt{Lupatotina Gas e Luce SRL} from a single
contributor. Concentration of two bills from one contributor is not a GDPR
violation per se, but it does mean the corpus is not balanced in a way that
distributes residual identification risk evenly across contributors; the same
applies to the four BC~Hydro bills, which trace to a single Stream A page. The
construct threat this raises is statistical, not legal, and it is recorded as such
in \S\ref{sec:validity}.

\subsection{Intellectual property}

The bill PDFs are the issuers' documents and may carry copyright in the layout,
typography and any printed graphics. Two licensing considerations apply. First,
fair-dealing for criticism, review and research in UK copyright law (Sections~30
and 29A of the Copyright, Designs and Patents Act 1988~\cite{cdpa1988}) covers the
research use here: the bills are processed as data for an experiment, not
republished as content, and the dissertation reproduces no bill image as a figure.
The dataset accompanying the source-code submission is consistent with the same
limited-research-use frame, and is not redistributed beyond the submission. Second,
the Anthropic and OpenAI terms of service in force during the April~2026 runs
permit the use of customer-supplied document content for the API call, and the
content is not retained by the providers for training under the no-training
default that both providers offer for paid API tiers; the project relied on this
default rather than on a separate enterprise data-processing agreement, and that
choice is recorded in \S\ref{sec:exec-repro} so a future replicator can confirm
the providers' policies have not changed.

The third licensing axis is on the project's own outputs. The source code is
released under an open-source licence (placed in \texttt{LICENSE} in the
repository) so a future replicator can extend it; the redacted dataset is not
released externally, given the GDPR considerations above; and the dissertation
itself remains the author's. None of these decisions is novel; they are the
default for an undergraduate project of this shape, and they are recorded here
because the BCS Code's clause on third-party legitimate rights is explicit about
intellectual property and unlawful disclosure of confidential information.

\section{Social and ethical considerations}
\label{sec:lsesp-social}

\subsection{Dual-use risk in template-free extraction}

The extraction pipeline studied here is general. It reads a bill and emits a
twelve-field structured record without any per-issuer rule, and the same code,
unmodified, would read a different document class. Three deployment patterns
follow from that property. The first is consumer-facing and benign: bill-comparison
services, energy-poverty triage tools, and accessibility wrappers that read bills
aloud or summarise them for users with low literacy in the bill language. The
second is enterprise back-office automation: accounts-payable processing, expense
reimbursement, and supplier-contract reconciliation. Both use the system to do
what the customer or the back-office user already wants to do.

The third pattern is the one that warrants attention. A pipeline that converts
bills to structured records at scale, with no template work and no upfront
investment, is exactly the kind of tool that lowers the cost of unauthorised
profiling. A landlord with a stack of incoming utility bills from prospective
tenants, an insurer with bills submitted as proof of address, or a debt buyer with
historical bills attached to delinquent accounts can extract consumption patterns
that the data subject did not consent to be analysed. The capability is not new;
specialist OCR vendors have offered it for years. The novelty is that template-free
extraction collapses the deployment cost, and an A++ ethical analysis has to take
that cost collapse seriously rather than treat it as someone else's problem.

The mitigations are not in the pipeline. They are in how a deployer of the
pipeline establishes the lawful basis for processing each batch of bills, retains
or discards the OCR text, and limits the downstream uses of the structured
records. The dissertation makes no claim about how a future deployer should
satisfy those obligations, but it does make a claim the literature does not make
clearly enough: that the difficulty of bill extraction has been overstated as a
barrier to misuse, and that the cost-vs-coverage results in
\S\ref{sec:rq2-tradeoff} should be read in that light. Bill extraction is now
cheap; the access controls are the security boundary, not the technical
difficulty.

\subsection{Bias, accessibility, and globalisation}

Three further axes of social impact apply.

\textbf{Bias and language coverage.} The mechanisms documented in
\S\ref{sec:rq1-modality} are not language-symmetric. M3 (locale-aware separator
handling on \texttt{consumption\_amount}) penalises European-locale numerics
specifically, and the four-bucket re-attribution in \S\ref{sec:failure-attribution}
shows that residual locale errors are LLM-side rather than OCR-side. A system
deployed in a multilingual market without the corrective in \S\ref{sec:design-norm}
would systematically under-extract on French, German, Italian and Spanish bills
relative to English bills. The bias is technical, not cultural, but its effect is
distributional: the populations less well served by the deployed system would be
those whose bills are written in the languages where the system is weakest. This
is the bias profile reported in the systematic-OCR-bias literature for non-Latin
scripts~\cite{nguyen2021post}, restated for European-locale numerics.

\textbf{Accessibility.} Bill comprehension is itself an accessibility problem.
Annex I of the EU Electricity Directive 2019/944~\cite{eu_directive_2019_944_annexI}
requires that bills be comprehensible to ordinary consumers, but field studies
consistently show that consumers struggle to identify the headline numbers on a
multi-page bill. A reliable template-free extractor is the kernel of an
assistive tool: a screen-reader that reads only the twelve relevant fields, a
plain-language summariser that turns a six-page bill into a three-line summary,
or a comparison view that aligns bill-to-bill changes for the same household. The
hallucination-vs-omission profile in \S\ref{sec:rq3-provider} matters for this
use case in a specific way: an assistive tool that hallucinates a bill total is
worse than one that misses a field, because the user has no independent way to
detect the wrong value. Sonnet's hallucination rate (twenty across the corpus,
against five for GPT-4o) is therefore an accessibility-relevant property, not
just a model-evaluation property.

\textbf{Globalisation.} The corpus covers four languages, all in Latin script
(\S\ref{sec:dataset-composition}). Non-Latin scripts (Arabic, Cyrillic, CJK) are
out of scope, and the per-language slice findings in \S\ref{sec:rq1-modality}
should not be extrapolated to language families not exercised. A globally
deployed system would need at minimum a separate Tesseract data file per
script for the OCR-mediated path, and would face vision-side performance changes
that the within-Latin results here cannot predict. This is the same caveat the
external-validity discussion in \S\ref{sec:validity} raises, restated as a global
deployment question rather than as a corpus-coverage question.

\section{Environmental and sustainability impact}
\label{sec:lsesp-environment}

The experiment produced sustainability-relevant measurements as a by-product.
The four-run corpus consumed 547,136 input tokens and 21,178 output tokens across
160 LLM calls, against a directly-measured monetary cost of \$1.79 in API charges
(Table~\ref{tab:cost}). Energy consumption per token is not published by either
provider, but published estimates put the order of magnitude at $10^{-3}$ Wh per
input token for current-generation hosted LLMs~\cite{patterson2021carbon}.
Applied to this run, the order-of-magnitude energy cost is on the order of $0.6$~kWh
total, which is below the energy cost of one hour of laptop CPU time for the OCR
pipeline.

The point is not the absolute number, which is small for a project at this
scale. It is the per-document scaling: a vision-mediated run on a corpus of
500{,}000 bills, the kind of corpus an enterprise back-office automation
deployment would handle, would cost on the order of several thousand kWh per
month at current per-token energy intensities. That figure is sensitive to the
two-page cap defended in \S\ref{sec:two-page-cap}: a hypothetical full-document
vision run on the same corpus, on average $3.9\times$ the input volume, scales
the energy cost proportionally. The cost-vs-coverage trade-off in
\S\ref{sec:rq2-tradeoff} is therefore not just a financial trade-off; it is also
a sustainability trade-off, and a deployer minimising carbon cost has direct
empirical reason to prefer the page-bounded vision pipeline over an unbounded
one when the additional pages do not actually carry required fields.

The OCR-mediated pipeline has a different sustainability profile. Tesseract is a
local computation on commodity CPU hardware. Its marginal carbon cost is the
electricity drawn by the host machine, and it does not require cloud-side
compute beyond the LLM call on the smaller text payload. On a corpus where every
required field already lives on the first two pages, the OCR-mediated path is
the lower-carbon option for both providers, by roughly the input-token-volume
ratio of \S\ref{sec:rq2-tradeoff}. The trade-off direction inverts on a corpus
that exercises later pages, where Tesseract latency scales with page count
(\S\ref{sec:two-page-cap}) but the underlying energy cost of one-pass page
rasterisation is still small relative to a vision-mode token bill. The
sustainability framing therefore reinforces the operational framing in
\S\ref{sec:rq2-tradeoff}: the cheapest extraction is the one that reads no more
of the document than the schema requires, regardless of modality.

\section{Software trustworthiness and computer security}
\label{sec:lsesp-trust}

A claim of correctness on a deployed system is a claim about more than its
average accuracy. The hallucination-vs-omission profile in
\S\ref{sec:rq3-provider} is the most consequential trustworthiness finding the
experiment produced: the two providers occupy distinct points in that space, and
a downstream system that treats one provider's outputs as interchangeable with
the other's is making an unstated assumption about which kind of failure it can
absorb. A bill-payment automation that ingests a hallucinated total is silently
wrong; one that ingests a missing total has the option of falling back to a
human reviewer. The provider choice is therefore a downstream-cost-tolerance
choice, and any deployment that does not document which kind of failure it can
absorb is opaque about what it has assumed.

The four-bucket re-attribution in \S\ref{sec:failure-attribution} is a related
trustworthiness contribution. The default OCR-vs-LLM heuristic of
\S\ref{sec:eval-framework} mis-attributes the cases discussed in
\S\ref{sec:diagnosis-heuristic}: contextual-inference failures, ground-truth
phrasing decisions, and locale-formatting on numerics. A diagnosis tool that
reports ``50\% OCR failure'' to its operator hides the fact that roughly half
of OCR-mode failures live in the LLM, the prompt or the annotation, and a
remediation programme based on that figure would invest in the wrong layer.
Truthful failure-attribution is part of software trustworthiness, and the
sharper four-bucket reading is the corrective.

The two security-adjacent threats are minor on this project's scope but
deserve naming. The first is prompt injection from bill content: a bill could
in principle contain text that overrides the extraction prompt, since the LLM
sees the bill content as part of the input. The risk is bounded here because
the prompt template treats bill content as data and not as instructions and the
output is parsed against a strict JSON schema (\S\ref{sec:exec-repro}), so a
malformed injection cannot escape the schema. The second is exfiltration of
bill content via the LLM call. Both providers' API terms in force at the time
of the runs commit to not using paid-tier customer content for training, and
the runs were configured under those defaults. A deployment under different
terms would need to re-evaluate this; the dissertation does not extrapolate
beyond the providers and tiers tested.

\section{Economic and commercial factors}
\label{sec:lsesp-economic}

The cost ranking in \S\ref{sec:rq2-tradeoff} (vision more expensive than
OCR, with the gap larger for the more expensive provider) places the per-bill
extraction cost at \$0.008--0.015 across the four conditions. That is below the
cost of one minute of human review time at any reasonable rate; on
extraction tasks where the schema is small and the per-bill review burden is
low, automated extraction is economically dominant for any corpus larger than a
few hundred bills. This is the headline commercial reason behind the recent
shift in the document-AI market away from bespoke per-issuer templates toward
LLM-mediated extraction, and the literature gap recorded in
\S\ref{sec:research-gap} is partly a market-following gap: comparison work
lagged the deployment trend.

The commercial implication of the four-bucket re-attribution
(\S\ref{sec:failure-attribution}) is more specific. A vendor whose roadmap
prioritises Tesseract replacement as the path to higher accuracy will
under-perform a vendor whose roadmap prioritises prompt design and
LLM-side-instruction hardening, on roughly half of the residual failures on
this document class, because half of those failures are not Tesseract-side.
The cost of getting that prioritisation wrong is not symmetric: OCR-engine
swaps are infrastructure work and prompt work is per-document-class work, and
the wrong investment programme can absorb several engineer-quarters before its
returns are evaluable.

A second commercial axis is per-issuer concentration. The corpus contains
thirty-three unique issuers, but four BC~Hydro bills, two~Affinity~Water bills,
two Vattenfall Europe Sales bills, and several other multi-document issuers
account for a large share of the per-issuer error mass
(\S\ref{sec:dataset-composition}). A commercial deployer reading the per-field
accuracy heatmap of \S\ref{sec:headline} should expect accuracy ceilings that
shift as the issuer mix shifts, and should not generalise the on-corpus
ceilings to a fresh issuer mix without re-evaluation. This is the same caveat
the external-validity discussion records, framed as a procurement question.

\section{Professional conduct and the BCS Code}
\label{sec:lsesp-bcs}

The BCS Code of Conduct~\cite{bcs2022} sets four obligations for members; each
applies to the project in a specific way.

\textbf{Public interest.} The Code's first clause requires due regard for
public health, privacy, security, the environment, and the legitimate rights of
third parties. The privacy obligations are addressed in \S\ref{sec:lsesp-legal},
the third-party rights in the intellectual-property paragraph, the security
obligations in \S\ref{sec:lsesp-trust}, and the environmental obligations in
\S\ref{sec:lsesp-environment}. The accessibility paragraph in
\S\ref{sec:lsesp-social} addresses the non-discrimination clause: the system's
language coverage limits would, if deployed without correction, distribute
benefit unevenly across language communities. The dual-use discussion in the
same section addresses the Code's clause on legitimate rights of third parties:
data subjects are third parties to a bill-extraction deployment, and the
deployer's obligations to them are not satisfied by the technical work alone.

\textbf{Professional competence and integrity.} The Code requires that
members claim only competence they can demonstrate, and present results
honestly. Two specific integrity moves are documented in this dissertation. The
falsification framing of \S\ref{sec:rq1-aggregate}--\S\ref{sec:rq3-provider}
fixes hypotheses before the data is seen, so a null result stays informative
against a pre-registered position; this is a stronger commitment than is
typical in undergraduate-project evaluations and was made specifically to
remove the temptation to re-write hypotheses after the fact. The four-bucket
re-attribution in \S\ref{sec:failure-attribution} is the second: rather than
report the strict heuristic's headline split as if it were ground truth, the
analysis names the heuristic's mis-attributions and supplies the corrective.
Both moves implement the integrity obligation rather than merely cite it.

\textbf{Duty to relevant authority.} The relevant authority is the supervisor
and the institution. Compliance is straightforward: the project follows the
KCL ethics process for human-subject considerations (none arise: no human
subjects participated, beyond the Stream~B contributors who consented to bill
release), the supervisor was consulted on the corpus-collection design, and
the dissertation accurately represents the experimental setup so a future
replicator could reconstruct it from the run directories.

\textbf{Duty to the profession.} The Code's fourth area covers reputation,
standards, and improvement of the profession. The most direct contribution
the dissertation makes is the diagnosis-heuristic critique: the substring
attribution rule is a widely used default, and the four-bucket re-classification
shows where it mis-attributes and what the corrective looks like. Sharing
that contribution through the open-source release of the analysis scripts
(\texttt{scripts/analysis/}) is the practical step required by this clause; a
future practitioner picking up the same heuristic on their own corpus has a
worked critique to start from.

\section{Summary}
\label{sec:lsesp-summary}

The legal status of utility bills is settled (GDPR Article~4(1), with lawful
bases for both corpus streams documented). The ethical surface is dual-use
rather than benign: template-free extraction lowers the deployment cost of
unauthorised profiling, and that cost collapse is a relevant fact about the
technology, not a deployment-level afterthought. Bias is locale-numeric and
not symmetric across languages. The sustainability framing reinforces the
operational framing: the cheapest extraction is the one that reads no more of
the document than the schema requires. The hallucination-versus-omission
profile is a software-trustworthiness fact, not just an accuracy fact, and the
four-bucket re-attribution is the integrity move the project's own diagnosis
heuristic asked of itself. None of the four BCS pillars carries a generic
restatement; each is connected to a specific section of the experimental
chapters, and the obligations the Code imposes on a member doing this kind of
work are not the same as the obligations on a member doing a different kind of
work.
